
\medskip
\noindent\SGARefTarget{sga3:xi:proposition:3:7}\textbf{Proposition 3.7.} ---
\textit{Let \(S\) be a locally noetherian prescheme, let \(I\) be a
directed partially ordered index set, let \((T_i)_{i\in I}\) be a
projective system of \(S\)-preschemes locally of finite type, let
\[
  T=\varprojlim_{i\in I}T_i
\]
be the projective-limit functor, let \(F\) be a functor over \(S\),
and let \(u:F\to T\) be an \(S\)-morphism.  Suppose that the following
conditions hold:}

\begin{enumerate}
\item[(a)] \SGARefTarget{sga3:xi:proposition:3:7:item:a}%
\textit{\(F\) is a sheaf for the faithfully flat quasi-compact
topology, commutes with inductive limits of rings, and commutes with
adic projective limits of local artinian rings.}

\item[(b)] \SGARefTarget{sga3:xi:proposition:3:7:item:b}%
\textit{The morphism \(u:F\to T\) is a monomorphism.}

\item[\(b'\)] \SGARefTarget{sga3:xi:proposition:3:7:item:b-prime}%
\textit{The morphism \(u:F\to T\) is infinitesimally étale.}

\item[(c)] \SGARefTarget{sga3:xi:proposition:3:7:item:c}%
\textit{For every point \(\xi\) of \(F\) with values in the spectrum
of a field \(k\), let
\(\xi_i\in T_i(\operatorname{Spec}k)\) be its image and let \(t_i\)
be the corresponding point of \(T_i\).  There exists \(i\in I\) such
that, for every \(j\geq i\), the transition morphism
\[
  p_{i,j}:T_j\longrightarrow T_i
\]
is étale at \(t_j\).}

\item[(d)] \SGARefTarget{sga3:xi:proposition:3:7:item:d}%
\textit{For every prescheme \(X\) locally of finite type over \(S\)
and every \(S\)-morphism \(X\to F\), the set of points \(x\in X\) at
which this morphism is infinitesimally étale is open.}
\end{enumerate}

\textit{Under these conditions, \(F\) is representable by a
prescheme locally of finite type over \(S\).}

Let us note at once that, in the case to be considered in the \SGARefLink{sga3:xi:section:4}{next
section}, conditions \SGARefLink{sga3:xi:proposition:3:7:item:c}{(c)} and \SGARefLink{sga3:xi:proposition:3:7:item:d}{(d)} will be verified by means of \SGARefLink{sga3:xi:corollary:3:8}{the
following corollary}.

\medskip
\noindent\SGARefTarget{sga3:xi:corollary:3:8}\textbf{Corollary 3.8.} ---
\textit{Assuming \SGARefLink{sga3:xi:proposition:3:7:item:a}{(a)}, \SGARefLink{sga3:xi:proposition:3:7:item:b}{(b)}, and \SGARefLink{sga3:xi:proposition:3:7:item:b-prime}{\((b')\)}, conditions \SGARefLink{sga3:xi:proposition:3:7:item:c}{(c)} and \SGARefLink{sga3:xi:proposition:3:7:item:d}{(d)} are
implied by the following conditions:}

\begin{enumerate}
\item[\(c'\)] \SGARefTarget{sga3:xi:corollary:3:8:item:c-prime}%
\textit{The \(T_i\) are smooth over \(S\), and the transition
morphisms \(p_{i,j}:T_j\to T_i\) are smooth.}

\item[\(d'\)] \SGARefTarget{sga3:xi:corollary:3:8:item:d-prime}%
\textit{For every point \(\xi\) of \(F\) with values in the spectrum
of a field \(k\), let \(t_i(\xi)\) be the point of \(T_i\) defined by
\(\xi\), let \(d_i(\xi)\) be the relative dimension of \(T_i\) over
\(S\) at \(t_i(\xi)\), and put
\[
  d(\xi)=\sup_i d_i(\xi).
\]
Then:}
  \begin{enumerate}
  \item[\(1^\circ\)] \SGARefTarget{sga3:xi:corollary:3:8:item:d-prime:1}%
  \textit{for every \(\xi\) as above,
  \(d(\xi)<+\infty\); and}

  \item[\(2^\circ\)] \SGARefTarget{sga3:xi:corollary:3:8:item:d-prime:2}%
  \textit{for every prescheme \(X\) locally of finite type over
  \(S\) and every \(S\)-monomorphism \(v:X\to F\), the function
  \[
    x\longmapsto d(\xi_x)
  \]
  on \(X\) is locally constant, where, for \(x\in X\), \(\xi_x\)
  denotes the point of \(F\) with values in \(\kappa(x)\) induced by
  \(v\).}
  \end{enumerate}
\end{enumerate}

\smallskip
\noindent\emph{Proof of Proposition
\SGARefLink{sga3:xi:proposition:3:7}{3.7}.}
Work under condition \SGARefLink{sga3:xi:proposition:3:7:item:c}{(c)}.  Let \(t\) be the element of
\(\operatorname{ens}(F)\) defined by \(\xi\) (cf.\ \SGARefLink{sga3:xi:remark:3:6}{3.6}), and put
\SGARefTarget{sga3:xi:display:3:7:local-ring}
\[
  \mathscr O_t=\varinjlim_i\mathscr O_{T_i,t_i}.
\]
Using condition \SGARefLink{sga3:xi:proposition:3:7:item:c}{(c)}, one sees readily that \(\mathscr O_t\) is a
noetherian local ring (EGA \(0_{\mathrm{IV}}\), 10.3.1.3).  Its
residue field \(\kappa(t)\) is the inductive limit of the residue
fields \(\kappa(t_i)\), and \(k\) is an extension of it.  If
\(\eta\) denotes \(\operatorname{Spec}\kappa(t)\), and \(\xi\)
denotes \(\operatorname{Spec}k\), there is a commutative diagram

\SGARefTarget{sga3:xi:diagram:3:7:residue-field-square}
\[
\begin{tikzcd}[column sep=huge,row sep=large]
  \xi \arrow[r] \arrow[d] & F \arrow[d] \\
  \eta \arrow[r] & T ,
\end{tikzcd}
\]
whose definition is evident and is left to the reader.  Since
\(\xi\to\eta\) is a covering for the fpqc topology, \(F\) is a sheaf
for that topology by \SGARefLink{sga3:xi:proposition:3:7:item:a}{(a)}, and \(F\to T\) is a monomorphism by \SGARefLink{sga3:xi:proposition:3:7:item:b}{(b)}, the
morphism \(\eta\to T\) factors uniquely through a morphism
\(\eta\to F\).  We now record \SGARefLink{sga3:xi:lemma:3:9}{the following lemma}.

\medskip
\noindent\SGARefTarget{sga3:xi:lemma:3:9}\textbf{Lemma 3.9.} ---
\textit{Under the hypotheses of
\SGARefLink{sga3:xi:proposition:3:7}{3.7}, let
\[
  T'=\operatorname{Spec}(A)
  \quad\text{and}\quad
  T'_0=\operatorname{Spec}(k(A))
\]
with \(T'\) a noetherian local scheme, and let
\(i:T'_0\to T'\) be the canonical immersion.  Suppose that a
commutative square of morphisms is given:}

\SGARefTarget{sga3:xi:diagram:3:9:local-square}
\[
\begin{tikzcd}[column sep=huge,row sep=large]
  T'_0 \arrow[r] \arrow[d,"i"'] & F \arrow[d] \\
  T' \arrow[r] & T .
\end{tikzcd}
\]
\textit{Then there exists a unique \(T\)-morphism \(T'\to F\).}

The proof is the proof of \(\SGARefLink{sga3:xi:proposition:3:1}{3.1}\), \((\mathrm{iii})\Rightarrow
(\mathrm{ii})\), in which the representability of the \(S\) occurring
in that statement was not used.  One uses the facts that \(F\) is a
sheaf for the fpqc topology, that it commutes with adic projective
limits of local artinian rings, here the rings
\[
  A/\operatorname{rad}(A)^{n+1},
\]
and that \(F\to T\) is both a monomorphism and infinitesimally étale.

We apply Lemma \SGARefLink{sga3:xi:lemma:3:9}{3.9} with
\[
  T'=\operatorname{Spec}(\mathscr O_t),
  \qquad
  T'_0=\operatorname{Spec}(\kappa(t)).
\]
We have just constructed \(T'_0\to F\), and the canonical
homomorphisms
\[
  \mathscr O_{T_i,t_i}\longrightarrow\mathscr O_t
\]
define a projective system of morphisms
\(\operatorname{Spec}(\mathscr O_t)\to T_i\), hence a canonical
morphism \(\operatorname{Spec}(\mathscr O_t)\to T\).
\SGARefTarget{sga3:xi:diagram:3:9:application-square}%
The corresponding square \(\mathrm{(x)}\) plainly commutes: by the
definition of \(T\), it is enough to check this after replacing \(T\)
by \(T_i\).  We therefore obtain a unique \(T\)-morphism

\SGARefTarget{sga3:xi:display:3:9:v-prime}
\[
  v':T'=\operatorname{Spec}(\mathscr O_t)\longrightarrow F.
\]
Since \(F\) commutes with inductive limits of rings, \SGARefLink{sga3:xi:display:3:9:v-prime}{this morphism}
factors, for sufficiently large \(i\), through

\SGARefTarget{sga3:xi:display:3:9:v-i-prime}
\[
  v'_i:T'_i=\operatorname{Spec}(\mathscr O_{T_i,t_i})
  \longrightarrow F.
\]
For such an \(i\), one has
\SGARefTarget{sga3:xi:display:3:9:local-isomorphism}
\[
  T'\xrightarrow{\sim}T'_i,
  \qquad\text{that is,}\qquad
  \mathscr O_{T_i,t_i}\xrightarrow{\sim}\mathscr O_t.
\]
Indeed, \(T'\to T'_i\) is faithfully flat and quasi-compact, hence an
effective epimorphism, so it is enough to show that it is a
monomorphism.  Suppose that two morphisms \(R\rightrightarrows T'\),
with \(R\) representable, have the same composite with
\(T'\to T'_i\).  Their composites with \(T'\to T\), which factors
through \(T'\to T'_i\), are then equal.  Hence their composites with
\(T'\to T_j\), and therefore with \(T'\to T'_j\), are equal for
every \(j\).  The two original morphisms are equal because \(T'\) is
the projective limit of the \(T'_j\) in the category
\((\mathbf{Sch})\).

Consider the diagram

\SGARefTarget{sga3:xi:diagram:3:9:factorization-triangle}
\[
\begin{tikzcd}[column sep=huge,row sep=large]
  T' \arrow[r,"v'"] \arrow[d] & F \arrow[d] \\
  T'_i \arrow[r] \arrow[ur,"v'_i"'] & T_i ,
\end{tikzcd}
\]
where the unlabeled arrows are the evident ones.  The square and the
upper triangle commute.  Since \(T'\to T'_i\) is an epimorphism, the
lower triangle also commutes.

The ring \(\mathscr O_{T_i,t_i}\) is the inductive limit of the rings
of affine open neighborhoods of \(t_i\) in \(T_i\).  Since \(F\)
commutes with inductive limits of rings, \(v'_i\) comes from a
morphism
\SGARefTarget{sga3:xi:display:3:7:v-t}
\[
  v_t:U_t\longrightarrow F
\]
defined on an open neighborhood \(U_t\) of \(t_i\) in \(T_i\).
After shrinking this neighborhood if necessary, \(v_t\) is a
\(T_i\)-morphism.  Indeed, \(T_i\), being locally of finite type over
\(S\), also commutes with inductive limits of rings.  The composite of
\(v_t\) with \(F\to T_i\) is a monomorphism, and therefore \(v_t\) is
a monomorphism.

We claim that \(v_t\) is infinitesimally étale at \(t_i\).  Since
\(F\to T\) is infinitesimally étale, this is equivalent to the
composite
\[
  T'_i\xrightarrow{v'_i}F\longrightarrow T
\]
being infinitesimally étale at \(t_i\), or to the induced morphism
\(T'_i\to T\) being infinitesimally étale.  Since
\(T'\xrightarrow{\sim}T'_i\), it is in turn equivalent to
\(T'\to T\) being infinitesimally étale.  This is immediate, because
that morphism is the projective limit of the infinitesimally étale
morphisms \(T'_j\to T_j\).

Apply condition \SGARefLink{sga3:xi:proposition:3:7:item:d}{(d)}, which has not yet been used.  It follows that
\(v_t\) is infinitesimally étale in a neighborhood of \(t_i\).
After replacing \(U_t\) by a smaller open, we may therefore suppose
that \(v_t\) is a monomorphism that is infinitesimally étale.

As \(t\) varies in \(\operatorname{ens}(F)\), the family of morphisms
\(v_t\) satisfies \SGARefLink{sga3:xi:proposition:3:5}{Proposition 3.5}, which gives the conclusion of
Proposition \SGARefLink{sga3:xi:proposition:3:7}{3.7}.

\smallskip
\noindent\emph{Proof of Corollary
\SGARefLink{sga3:xi:corollary:3:8}{3.8}.}
Assume conditions \SGARefLink{sga3:xi:corollary:3:8:item:c-prime}{\((c')\)} and \SGARefLink{sga3:xi:corollary:3:8:item:d-prime}{\((d')\)}.  With the notation of
\SGARefLink{sga3:xi:corollary:3:8:item:d-prime}{\((d')\)}, \(d_i(\xi)\) is constant for all sufficiently large \(i\).
Consequently, the relative dimension of \(T_j\) over \(T_i\) at
\(t_j\), which is
\[
  d_j(\xi)-d_i(\xi),
\]
is zero.  Thus \(T_j\to T_i\) is étale at \(t_j\), which proves
condition \SGARefLink{sga3:xi:proposition:3:7:item:c}{(c)}.

The preceding proof now gives, for each
\(t\in\operatorname{ens}(F)\), an index \(i\), an open neighborhood
\(U_t\) of \(t_i\), and a morphism \(U_t\to F\) that is a
monomorphism and is infinitesimally étale at \(t_i\).  It remains to
show that \SGARefLink{sga3:xi:display:3:7:v-t}{this morphism} is infinitesimally étale throughout a
neighborhood of \(t_i\).  In condition \SGARefLink{sga3:xi:corollary:3:8:item:d-prime}{\((d')\)}, \SGARefLink{sga3:xi:corollary:3:8:item:d-prime:2}{\(2^\circ\)}, take
\(X=U_t\).  Replacing \(U_t\), if necessary, by the connected
component of \(t_i\) in \(U_t\), we may suppose that

\SGARefTarget{sga3:xi:display:3:8:dimension-constant}
\[
  d(\xi_x)=d(\xi_{t_i})
  \qquad\text{for every }x\in U_t.
\]
Now \(d(\xi_{t_i})\) is the relative dimension of \(T_i\) over \(S\)
at \(t_i\).  Since that relative dimension remains constant on
\(U_t\), we also have, for every \(x\in U_t\),

\SGARefTarget{sga3:xi:display:3:8:relative-dimension-star}
\[
  \tag{*}
  d(\xi_x)=\text{the relative dimension of \(U_t\) over \(S\) at \(x\)}.
\]

On the other hand, using the notation of the proof of
\SGARefLink{sga3:xi:proposition:3:7}{3.7}, consider any local
noetherian prescheme
\[
  R'=\operatorname{Spec}(A),
  \qquad
  R'_0=\operatorname{Spec}(k(A)).
\]
Every morphism \(R'\to F\) whose induced point \(R'_0\to F\)
represents \(t\in\operatorname{ens}(F)\), that is, sends the closed
point of \(R'\) to \(t\), factors uniquely through \(T'\).  The proof
is the proof of \SGARefLink{sga3:xi:proposition:3:1}{3.1} or of Lemma
\SGARefLink{sga3:xi:lemma:3:9}{3.9}: one uses that \(T'\to F\) is a
monomorphism infinitesimally étale at \(t\), and that \(T'\) is an
fpqc sheaf commuting with adic projective limits.

Apply this result to the morphism
\[
  \operatorname{Spec}(\mathscr O_{U_t,x})=R'\longrightarrow F
\]
induced by \(v_t\), and to the point
\(t_x\in\operatorname{ens}(F)\) that is the image of the closed point
of \(R'\), or equivalently the image of \(x\) under \(v_t\).  We
obtain a factorization

\SGARefTarget{sga3:xi:display:3:8:local-factorization}
\[
  \operatorname{Spec}(\mathscr O_{U_t,x})
  \longrightarrow
  \operatorname{Spec}(\mathscr O_{t_x})
  \longrightarrow F.
\]
The second arrow is infinitesimally étale at \(t_x\).  To prove that
the composite is infinitesimally étale at \(x\), it is therefore
enough to prove the same for the first arrow at \(x\).

By \SGARefLink{sga3:xi:diagram:3:9:application-square}{\(\mathrm{(x)}\)}
above and formula
\SGARefLink{sga3:xi:display:3:8:relative-dimension-star}{\((*)\)},
the first arrow is induced by a local \(S\)-homomorphism between the
localizations, at \(x\) and \(t_x\), of the smooth \(S\)-preschemes
\(U_t\) and \(U_{t_x}\), which have the same relative dimension
\(d\) over \(S\) at those points.  By EGA I, 6.5.1, this local
homomorphism is induced by an \(S\)-morphism

\SGARefTarget{sga3:xi:display:3:8:w}
\[
  w:U\longrightarrow V,
\]
where \(U\) is an open neighborhood of \(x\) in \(U_t\) and
\(V=U_{t_x}\).  It remains to prove that \(w\) is étale at \(x\).
Moreover, \(U\) and \(V\) carry monomorphisms to \(F\), and, after
shrinking \(U\) once more, \(w\) is an \(F\)-morphism.  Hence \(w\)
is a monomorphism.  It is now enough to prove \SGARefLink{sga3:xi:lemma:3:10}{the following lemma}.

\medskip
\noindent\SGARefTarget{sga3:xi:lemma:3:10}\textbf{Lemma 3.10.} ---
\textit{Let \(U\) and \(V\) be two smooth \(S\)-preschemes of the
same relative dimension \(d\) over \(S\), and let \(w:U\to V\) be an
\(S\)-morphism that is a monomorphism.  Then \(w\) is an open
immersion, and hence is étale.}

By SGA 1, I, 5.7, we reduce to the case in which \(S\) is the
spectrum of a field, which may be assumed algebraically closed.  By
SGA 1, I, 5.1, it is enough to prove that \(w\) is étale, and it is
enough to do so at the closed points of \(U\).  Let \(x\) be such a
point and put \(y=w(x)\).  Choose a regular system of parameters
\(f_1,\ldots,f_d\) of \(\mathscr O_{V,y}\).  Then
\[
  A=
  \mathscr O_{U,x}\big/\sum_i f_i\mathscr O_{U,x}
\]
is the trivial extension of
\[
  k=
  \mathscr O_{V,y}\big/\sum_i f_i\mathscr O_{V,y};
\]
that is, \(A\simeq k\).  Indeed, since \(w\) is a monomorphism, so is
the structural morphism
\[
  \operatorname{Spec}(A)\longrightarrow\operatorname{Spec}(k)
\]
obtained from it by base change.  Since \(\mathscr O_{U,x}\) is a
regular local ring of dimension \(d\), the \(f_i\) form a regular
system of parameters of that ring.  It follows immediately that
\(w\) is étale at \(x\).  This proves Lemma
\SGARefLink{sga3:xi:lemma:3:10}{3.10}, and hence Corollary
\SGARefLink{sga3:xi:corollary:3:8}{3.8}.

\medskip
\noindent\SGARefTarget{sga3:xi:corollary:3:11}\textbf{Corollary 3.11.} ---
\textit{Under the hypotheses of
\SGARefLink{sga3:xi:corollary:3:8}{3.8}, for every quasi-compact open
subfunctor \(U\) of \(F\) that is separated over \(S\), there exists
\(i\in I\) such that, for every \(j\geq i\), the morphism
\[
  u_j|_U:U\longrightarrow T_j
\]
is an open immersion.  In particular, if the \(T_i\) are quasi-affine
over \(S\), then every open subfunctor \(U\) of \(F\) that is
quasi-compact over \(S\), that is, of finite type over \(S\), is
quasi-affine over \(S\).}

The proof of Proposition
\SGARefLink{sga3:xi:proposition:3:7}{3.7} shows that for every
\(t\in\operatorname{ens}(F)\) there exists \(i\in I\) such that
\(u_i:F\to T_i\) is a local isomorphism at \(t\), and then \(u_j\)
is a local isomorphism at that point for every \(j\geq i\).  By
quasi-compactness, \(i\) may be chosen independently of \(x\in U\).
It remains to prove that, for sufficiently large \(i\),
\[
  u_i|_U:U\longrightarrow T_i
\]
is a monomorphism.

Since \(U\to T\) is a monomorphism, the intersection of the
equivalence relations
\[
  U\times_{T_i}U\ \subset\ U\times_S U
\]
is the diagonal.  The prescheme \(U\times_S U\) is noetherian, and
the \(U\times_{T_i}U\) are closed subpreschemes.  It follows that one
of them is already the diagonal; equivalently, \(u_i|_U\) is a
monomorphism.  This proves the first assertion of Corollary
\SGARefLink{sga3:xi:corollary:3:11}{3.11}, and the second follows
immediately.

